\documentclass[11pt,a4paper]{article}
\usepackage[utf8]{inputenc}
\usepackage[T1]{fontenc}
\usepackage{geometry}
\usepackage{graphicx}
\usepackage{booktabs}
\usepackage{array}
\usepackage{float}
\usepackage{amsmath}
\usepackage{amssymb}
\usepackage{hyperref}
\usepackage{pgfplots}
\usepackage{csquotes}
\usepackage[backend=biber,style=numeric-comp,sorting=none,natbib=true]{biblatex}
\addbibresource{reference.bib}

\geometry{margin=1in}
\pgfplotsset{compat=1.18}

\hypersetup{
  colorlinks=true,
  linkcolor=black,
  urlcolor=black,
  citecolor=black,
  pdftitle={Week 5--6 Progress Report: Phase 3},
  pdfauthor={Deep Shukla}
}

\title{Week 5--6 Progress Report\\Phase 3: Evaluation and Results}
\author{Deep Shukla}
\date{April 20, 2026}

\begin{document}
\maketitle

\section{Introduction}

This phase focuses on evaluating the implemented dataset retrieval prototype for Swiss Open Government Data (OGD). The system integrates metadata retrieval from opendata.swiss, query processing, feature engineering, fuzzy ranking, baseline comparison, and explainable outputs. The overall artifact development follows a Design Science Research orientation, where implementation and evaluation are carried out iteratively \citep{hevner2004,peffers2007,thesisPhase2_2026}.

The core retrieval challenge in OGD portals is that discoverability depends heavily on metadata quality and query vagueness. To address this, the implemented system combines keyword similarity with quality-aware fuzzy scoring and interpretable explanations \citep{janssen2012,attard2015,zadeh1965,mamdani1974,doshi2017}.

\section{System Architecture}

The evaluated artifact preserves the modular architecture specified in the previous phase, with clear separation between query processing, candidate retrieval, feature extraction, fuzzy inference, and ranking output \citep{thesisPhase2_2026}.

\subsection{Architecture Diagram}

The implemented component-level architecture is summarized in Table~\ref{tab:modules}.

\begin{table}[H]
\centering
\begin{tabular}{p{0.34\textwidth} p{0.56\textwidth}}
\toprule
\textbf{Module} & \textbf{Role} \\
\midrule
\texttt{code/prototype/app.py} & Prototype entrypoint wrapper. \\
\texttt{code/prototype/swiss\_ogd\_portal.py} & Main interactive retrieval application. \\
\texttt{code/data\_collection/ckan\_api\_client.py} & CKAN metadata acquisition layer. \\
\texttt{code/query\_processing/query\_parser.py} & Query parsing and intent cue extraction. \\
\texttt{code/ranking/baseline\_keyword.py} & BM25/keyword baseline retrieval. \\
\texttt{code/ranking/fuzzy\_ranker.py} & Fuzzy ranking orchestration. \\
\texttt{code/ranking/ai\_semantic\_baseline.py} & Optional semantic baseline. \\
\texttt{evaluation/evaluation\_framework.py} & IR metric computation and evaluation logic. \\
\texttt{evaluation/experiment\_runner.py} & End-to-end experiment execution. \\
\bottomrule
\end{tabular}
\caption{Implemented prototype modules used in Phase 3 evaluation}
\label{tab:modules}
\end{table}

\subsection{Pipeline Description}

The runtime pipeline is:
\begin{enumerate}
  \item Parse and normalize query.
  \item Retrieve candidate datasets via CKAN/local snapshot.
  \item Compute feature signals.
  \item Apply ranking method (BM25, fuzzy, optional SBERT).
  \item Produce top-$k$ ranked list and explanations.
  \item Evaluate rankings with graded relevance metrics.
\end{enumerate}

This pipeline aligns with IR evaluation practice and explainable decision support requirements \citep{manning2008,jarvelin2002}.

\section{Data Acquisition and Processing}

\subsection{Preprocessing}

The evaluation uses the processed metadata snapshot in \texttt{data/raw/ogd\_metadata\_20260306\_183841.json}, sourced from opendata.swiss metadata exports \citep{opendataswiss2026,ogdSnapshot2026}. Records are normalized into retrieval-ready fields (title, description, keywords, themes, resource metadata, recency indicators, and completeness attributes).

\begin{table}[H]
\centering
\begin{tabular}{p{0.45\textwidth} p{0.24\textwidth}}
\toprule
\textbf{Statistic} & \textbf{Value} \\
\midrule
Total datasets & 500 \\
Unique organizations & 3 \\
Average resources per dataset & 6.84 \\
Median resources per dataset & 9.0 \\
Average keywords per dataset & 8.65 \\
Median keywords per dataset & 7.0 \\
Average completeness score & 0.75 \\
Creation range & 2022-02-28 to 2026-03-05 \\
Modification range & 2026-03-06 to 2026-03-06 \\
\bottomrule
\end{tabular}
\caption{Dataset snapshot profile}
\label{tab:snapshot}
\end{table}

\subsection{Query Processing}

The query set is stored in \texttt{evaluation/benchmark\_queries\_v2.json} and includes 15 benchmark queries across environment, mobility, and cross-domain tasks \citep{benchmarkQueries2026}. The parser extracts keywords, language cues, and intent modifiers.

\begin{table}[H]
\centering
\begin{tabular}{p{0.12\textwidth} p{0.14\textwidth} p{0.42\textwidth} p{0.22\textwidth}}
\toprule
\textbf{ID} & \textbf{Domain} & \textbf{Query} & \textbf{Intent} \\
\midrule
ENV-01 & environment & Aktuelle Luftqualit\"atsdaten f\"ur Schweizer St\"adte & Air quality retrieval \\
ENV-03 & environment & Wasserqualit\"at in Schweizer Seen und Fl\"ussen & Water quality retrieval \\
MOB-01 & mobility & Oeffentlicher Verkehr Fahrplandaten Schweiz & Transit schedule data \\
MOB-03 & mobility & Elektromobilit\"at Ladestationen Schweiz & EV charging data \\
MOB-05 & mobility & Parkpl\"atze und Parkh\"auser Echtzeitdaten & Parking availability \\
XD-01 & cross & Bev\"olkerungsentwicklung und Demografie Schweiz & Demographic statistics \\
XD-02 & cross & Energieverbrauch nach Kantonen & Energy by canton \\
XD-05 & cross & Wirtschaftsstatistik und BIP nach Region & Regional economy \\
\bottomrule
\end{tabular}
\caption{Representative benchmark queries}
\label{tab:queries}
\end{table}

\section{Feature Engineering}

The fuzzy model uses four primary input criteria per query-dataset pair.

\subsection{Keyword Similarity}

Similarity is computed using keyword/BM25-style lexical matching for baseline retrieval and as an input signal for fuzzy ranking \citep{robertson2009,manning2008}.

\subsection{Recency}

Recency is derived from metadata modification timestamps. More recent datasets receive stronger recency contribution.

\subsection{Completeness}

Completeness is estimated from metadata field coverage (e.g., title, description, tags, resources, license/organization fields), normalized to $[0,1]$.

\subsection{Category Match}

Category match reflects thematic overlap between query intent and dataset category/theme indicators.

\section{Fuzzy Model Design}

\subsection{Fuzzy Variables}

The fuzzy system combines similarity, recency, completeness, and category/metadata-quality cues into a final relevance output.

\subsection{Fallback Strategy for Incomplete Rule Coverage}

When full rule activation is weak or unavailable, the prototype includes fallback ranking behavior to ensure robust output generation in all query scenarios.

\subsection{Rationale}

This design keeps ranking interpretable while allowing gradual relevance reasoning, which is central to the thesis objective \citep{zadeh1965,mamdani1974,doshi2017}.

\section{Membership Functions}

\subsection{Triangular Membership Function}

Triangular memberships are used for computational simplicity and interpretability.

\subsection{Parameter Selection}

Parameters follow prototype calibration and practical distribution ranges observed in metadata features.

\section{Fuzzy Rule Base}

\subsection{Primary Rules (High Similarity)}

Rules prioritize high lexical/semantic alignment together with good metadata quality.

\subsection{Secondary Rules (Moderate Similarity)}

Rules support exploratory retrieval by rewarding moderate similarity when quality signals are strong.

\subsection{Penalty Rules (Low Quality Conditions)}

Rules reduce ranking strength under sparse/outdated or weakly aligned metadata conditions.

\section{Fuzzy Inference Mechanism}

\subsection{Inference Steps}

The prototype applies fuzzification, rule evaluation, aggregation, and defuzzification to generate final ranking scores.

\section{Defuzzification}

The crisp relevance score is produced after output aggregation and used for final sorting in fuzzy retrieval lists.

\section{Ranking Mechanism}

The evaluated retrieval methods are:
\begin{enumerate}
  \item BM25/keyword baseline,
  \item fuzzy ranking (main system),
  \item optional SBERT/semantic baseline.
\end{enumerate}

Each method returns top-10 results per query for evaluation.

\section{Explainability}

For fuzzy results, explanations are generated from feature values. A typical explanation is: ``Dataset ranked high because it strongly matches query terms, is recent, and has complete metadata.'' This supports transparent decision rationale in contrast to black-box retrieval behavior \citep{doshi2017}.

\section{Prototype Implementation}

\subsection{Implementation Environment}

The prototype is implemented in Python with modular components and a Streamlit interface \citep{streamlit2026}.

\subsection{Data Integration}

Data integration is performed through CKAN API access and local metadata snapshots \citep{ckan2026,opendataswiss2026,ogdSnapshot2026}.

\subsection{Feature Computation Pipeline}

Feature computation includes similarity, recency, completeness, and category/theme signals for each candidate dataset.

\subsection{Fuzzy Inference Integration}

The fuzzy engine is integrated as a scoring layer in the ranking pipeline.

\subsection{Ranking and Output Generation}

Per query, each method returns a ranked top-10 list with method-specific scores and metadata.

\subsection{System Integration}

Integration testing confirms that module imports, parser behavior, baseline retrieval, fuzzy scoring, explanation generation, and configuration handling are operational \citep{integrationTest2026}.

\section{Preliminary Validation}

\subsection{Results}

The test input $\{\text{recency}=10,\text{completeness}=0.85,\text{thematic\_similarity}=0.75,\text{resource\_availability}=5\}$ produced a fuzzy relevance score of 91.5/100 in integration output \citep{integrationTest2026}.

Aggregate retrieval metrics over 15 queries are shown in Table~\ref{tab:aggregate}.

\begin{table}[H]
\centering
\begin{tabular}{p{0.22\textwidth} p{0.09\textwidth} p{0.09\textwidth} p{0.10\textwidth} p{0.11\textwidth} p{0.11\textwidth} p{0.08\textwidth}}
\toprule
\textbf{Method} & \textbf{MAP} & \textbf{P@10} & \textbf{R@10} & \textbf{nDCG@10} & \textbf{nDCG@20} & \textbf{MRR} \\
\midrule
portal\_default & 0.9877 & 0.9933 & 0.3374 & 0.8129 & 0.8489 & 1.0000 \\
keyword\_bm25 & 0.7992 & 0.8667 & 0.2955 & 0.7393 & 0.7370 & 0.9667 \\
metadata\_quality & 0.9824 & 0.9933 & 0.3374 & 0.7631 & 0.8233 & 1.0000 \\
fuzzy\_hcir & 0.8367 & 0.7133 & 0.2430 & 0.5796 & 0.7305 & 0.8800 \\
\bottomrule
\end{tabular}
\caption{Aggregate performance comparison}
\label{tab:aggregate}
\end{table}

\subsection{Interpretability}

Fuzzy outputs provide explicit feature-level traces and text explanations, supporting interpretability requirements for public-sector retrieval settings.

\subsection{Limitations}

The current snapshot shows stronger performance of portal default ranking in top-rank metrics. Fuzzy behavior indicates comparatively larger gain from nDCG@10 to nDCG@20, suggesting scope for top-rank tuning.

\section{Limitations}

Two key boundaries remain:
\begin{enumerate}
  \item Human rater data are not yet populated in the current repository snapshot.
  \item Semantic baseline quality can be further refined.
\end{enumerate}

\section{Evaluation Plan}

\subsection{Evaluation Objectives}

The evaluation targets comparison of fuzzy ranking against baseline methods and assessment of interpretability effects.

\subsection{Dataset and Query Set}

The evaluation uses:
\begin{itemize}
  \item 500-dataset metadata snapshot,
  \item 15 benchmark queries,
  \item 450 graded relevance judgments.
\end{itemize}

\subsection{Baseline Methods}

BM25 baseline and optional semantic retrieval are used as comparators to fuzzy ranking.

\subsection{Evaluation Metrics}

Metrics include MAP, P@K, R@K, nDCG, MRR, and F1@K \citep{manning2008,jarvelin2002}.

\subsubsection{Normalized Discounted Cumulative Gain (nDCG)}

\begin{equation}
\mathrm{DCG}@K = \sum_{i=1}^{K} \frac{2^{rel_i}-1}{\log_2(i+1)}
\end{equation}

\begin{equation}
\mathrm{nDCG}@K = \frac{\mathrm{DCG}@K}{\mathrm{IDCG}@K}
\end{equation}

\subsection{Interpretability Evaluation}

\subsubsection{User Study Design}

Planned user evaluation includes 15--20 participants, within-subject comparison, graded judgments, and perceived trust/usability assessment.

\subsubsection{Evaluation Procedure}

Inter-rater agreement will be assessed with Cohen/Fleiss statistics, and usability can be assessed with SUS-style measures \citep{cohen1960,fleiss1971,brooke1996}.

\subsection{Experimental Procedure}

The implemented experiment flow is:
\begin{enumerate}
  \item load query set,
  \item retrieve top-10 per method,
  \item assign/consume graded relevance labels,
  \item compute per-query metrics,
  \item aggregate and export outputs.
\end{enumerate}

Generated evaluation artifacts include results JSON, relevance labels, metrics tables, summary text, and comparative plots \citep{experimentResults2026,groundTruthAuto2026,experimentRunnerCode2026}.

\subsection{Expected Outcomes}

Expected outcomes are improved interpretability with competitive retrieval quality after fuzzy parameter refinement and stronger evidence from human-rated judgments.

\section{Conclusion}

Phase 3 delivers an end-to-end evaluated prototype with modular retrieval methods, fuzzy explainability, quantitative metric outputs, and reproducible evaluation artifacts. The document format and structure now follow the same academic pattern used in Phase 2, while incorporating the full Phase 3 evaluation content.

\subsection{Next Phase of Work}

The next phase will prioritize: (i) human-rater study execution, (ii) top-rank fuzzy calibration, and (iii) final thesis consolidation of quantitative and interpretability findings.

\printbibliography

\end{document}
